\insimgspl{forprint/002_folio_01_recto.jpg}
%\vskip -10pt

%\Rline{\krnt, folio 1\,\RE.}
\smallskip

\insimgspl{forprint/003_folio_01_verso.jpg}
%\vskip -10pt

%\Rline{\krnt, folio 1\,\VE.}
\smallskip

\insimgspl{forprint/004_folio_02_recto.jpg}
\vskip -10pt

\Rline{\krnt, folio 1\,\RE, 1\,\VE, 2\,\RE.}


\newpage
~\vskip 1cm

\begin{sans}
\begin{center}
{\Huge\bfseries \color{deepmaroon}{koraṇṭakapaddhatiḥ}}
\bigskip

{\Large\bfseries śrīmahāgaṇapataye namaḥ\!॥ 
śrīgurubhyo namaḥ\!॥\endnote{The marginal note in the top left corner says, vajrolīyoga\-paddhati\!\textsanskrit{॥} 
kuraṇṭakapaddhati, voralārtha (?), yogaśāstrārtha.}}
\end{center}
\bigskip

saṃsārataptānāṃ\endnote{saṃsāra°] P, saṃsārāsāra° M} % 
atyantaviṣayāsaktānāṃ straiṇānāṃ jāti\-bhraṣṭāṇāṃ atisāhasakarmaकर्तॄणां,  tatkṛte iyaṃ kapālakuraṃṭakakṛtahaṭhābhyāsapaddhatiḥ\!।\endnote{°paddhatiḥ] \textit{emend.}, °paddhaty° M} tadgata\-padā\-rthāḥ\endnote{tadgata°] \textit{emend.}, aṃtargata° M} sādhanakarmāṇi\endnote{°karmāṇi] P, °karmāṇī M} ca  likhyante\!।

maṭhikālakṣaṇam\!॥\\ % section heading
caturhastapramāṇamūrdhvaṃ\endnote{ūrdhvaṃ] P, ūrdhva° M} tiryak\!। bhasmamaṭhikā mudrābhyāsārtham\!। āraktamṛttikāmaṭhikā āsanābhyā\-sārtham\!। sudhāmaṭhikā bastyādyabhyāsārtham\endnote{bastyādya°] \textit{conj.}, bastyārthā° M}\!। vyāghrādicarmamaṭhikā śayanārtham\!। tūlavastrādimaṭhikā vajrolyartham\!। tridhanuṣyordhvamekadhanu\-ṣya\-tirya\-ksvargādyāsanārtham\!। samīpaṃ\endnote{samīpaṃ] P, samīpa M} gūḍhaṃ nirvedhaṃ malādityāgārthaṃ sthalam\!। mānuṣarahitajalakulyā snānārtham\!।

yamaniyamalakṣaṇam\!॥\\ % section heading
ahiṃsā satyamasteyamasaṅgo hrīrasañcayaḥ\!।\\
āstikyaṃ brahmacaryaṃ ca maunaṃ\endnote{maunaṃ] P, maunyaṃ M} sthairyaṃ kṣamābhayam\!॥\\
śaucaṃ japastapo homaḥ śra\insmark{0pt}{4pt}{ddhā}tithyaṃ madarcanam\!।\\ %[folio 1v ends at śraddhā]
tīrthāṭanaṃ parārthehā tuṣṭirācāryasevanam\!॥

kāyavāṅmanobhiḥ pareṣāṃ duḥkhadānābhāvaḥ \textbf{ahiṃsā}\!। yathārthaṃ\endnote{yathārthaṃ] P, yāthārthyaṃ M}
ghātāprāpakaṃ anudvegakaraṃ hitaṃ bhāṣaṇaṃ \textbf{satyam}\!। śarīramanovāgbhiḥ paravastvabhilāṣābhāvaḥ \textbf{asteyam}\!। abhyāsapratiyogijanānāmasnehaḥ  \textbf{asaṅgaḥ}\!। abhyāsanāśakeṣu dharmādharmeṣu lajjā \textbf{hrīḥ}\!। abhyāsopayogivastvanyavastusaṅgrāhābhāvaḥ\endnote{°grāhā°] P, ++ (damaged) M} \textbf{asañcayaḥ}\!। matsyendragorakṣakākikāpālikādīnāṃ śāstreṣu viśvāsaḥ \textbf{āstikyam}\!। upasthadvārā vīryabindorapatanaṃ raja[-]ādyākarṣaṇaṃ ca \textbf{brahmacaryam}\!। grāmyaviṣayakathāvivarjanaṃ \textbf{maunam}\!। dāraputrādipālanācintayā\endnote{dāra°] \textit{emend.}, dārā° M. °ācintayā] \textit{conj.}: °acintayā M} manaso 'cāñcalyaṃ \textbf{sthairyam}\!। śārīramānasacorarājaduṣṭajanaduḥkhe teṣāṃ daṇḍakaraṇasāmarthye sati tatsahanaṃ \textbf{kṣamā}\!। ahaṃkārāspadaśarīrādi mamatvādisthānaputragṛhādi anityam,  ātmā nityaḥ saccidānandaśceti jñānena nirbhayatvam \textbf{a\-bhayam}\!।

bastidhautigajakaraṇādinā\endnote{°karaṇādinā] P, karṇādinā M} śarīraśuddhiḥ, bhagavannāmasaṃkīrtanena vākśuddhiḥ, 
bhaga\insmark{0pt}{0pt}{va}ddhyānena\endnote{bhagavad°] P, bhagava° M} mānasaśuddhiḥ \textbf{śaucam}\!।

upāsyadevatāmantrasya nityāvartanaṃ \textbf{japaḥ}\!। svadharmācaraṇaṃ \textbf{tapaḥ}\!। nityahomaḥ gṛhasthāśramiṇāṃ brahmacāriṇāṃ vānaprasthānāṃ ca\endnote{ca] P, \textit{omitted.} M}, yatīnāṃ brahmabandhustryādīnāṃ\endnote{brahma°] P, braṃha° M} vaiṣṇavatantrādibhagavatpūjāṅgahomaḥ \textbf{homaḥ}\!। śāstraśrutavidhikaraṇautsukyaṃ \textbf{śraddhā}\!। svamārgahaṭhādiyuktyabhyastānāṃ satkāraḥ \textbf{ātithyam}\!। sarvabhūteṣu dayā, yena kena santuṣṭiḥ, sarvatrātmānudarśanaṃ bhagavatpūjā\!।  abhyāsasiddhānāṃ darśanārthaṃ deśe deśe  paryaṭanaṃ \textbf{tīrthāṭanam}\!। śiṣyābhyāsasiddhyarthaṃ\endnote{°siddhy°] P, °sidhy° M} yatnaḥ \textbf{parārthehā}\!। yadṛcchayā prāptena santoṣaḥ \textbf{tuṣṭiḥ}\!। śāstraniṣṇātānubhavaniṣṇātaguroḥ anugrahasaṃpādanam \textbf{ācāryasevanam}\!।
\medskip

[āsanāni\!॥]\\  
atha ṣaṭkarmayogyatāpratipādanāya āsanāni likhyante\!।
\end{sans}

\newpage
~\vskip 1cm
\begin{center}
\begin{eng}
\booktitle{Koraṇṭaka’s Manual [of Haṭhayoga]}
\medskip

{\large\bfseries Homage to the venerable and mighty Gaṇapati.\\[4pt] 
Homage to the venerable Gurus.}
\end{eng}
\end{center}
\medskip

\begin{multicols}{2}
\begin{eng}
This is a manual on the practice of Haṭha[yoga] composed by Kapālakuraṇṭaka for the sake of those afflicted by the pain of transmigration, those excessively attached to sense objects, those obsessed with women, those fallen from caste, and [even] those who perform the most egregious actions. The topics in it and the techniques of the practice shall be written [here].
\bigskip

\Section[Characteristics of the Yogi’s Hut]{Characteristics of the [Yogi’s] Hut}

The measurement is four forearm lengths high and wide. For the practice of seals (\textit{mudrā}), a hut [daubed] with ash [should be used]. For the practice of postures (\textit{āsana}), a hut [made] of reddish soil. For the practice of the [six cleansing techniques,] such as the enema (\textit{basti}), a hut [daubed] with plaster. For lying down, the hut should have an animal skin, such as that of a tiger. For \textit{vajroli}, it should have a cotton cloth and the like. For the practice of postures [requiring a rope], such as the posture [of Climbing up to] Heaven, it should [measure] three bow lengths high and one bow length wide. For excreting faeces and so on, there should be a place that is nearby, hidden and closed off. For bathing, there should be a stream of water free of people.
\bigskip

\Section{Characteristics of the Behavioural Codes (\textit{yama}) and Observances (\textit{niyama})}

Non-violence, truthfulness, not stealing, non-attachment, shame, not accumulating [goods], belief, celibacy, silence, steadiness, patience and fearlessness [are the  \textit{yama}s]. Purity, reciting mantras, asceticism, [performing] a fire \mbox{sacrifice,} confidence, hospitality, worship of me (i.e., god), travelling to sacred places, striving to help others, satisfaction and serving the teacher [are the \textit{niyama}s].\endnote{\hspace*{.01cm} Both verses were borrowed from the \textit{Bhāgavata\-purāṇa} (11.19.33–34).  In their original context, the compound \textit{madarcana} means ‘worship of me (i.e., Kṛṣṇa).’}

\textbf{Non-violence} is not causing suffering to others with one’s body, speech or thought. \textbf{Truthfulness} is speech that is factual, leads to no harm, causes no agitation and is beneficial. \textbf{Not stealing} is the absence of desire for the possessions of others in body, mind and speech. \textbf{Non-attachment} is the absence of affection towards people who oppose the practice [of yoga]. \textbf{Shame} is compunction for righteous and unrighteous [actions] that destroy one’s practice [of yoga]. \textbf{Not accumulating} is not collecting goods that are different to those needed for the practice [of yoga]. \textbf{Belief} is confidence in the scriptures of Matsyendra, Gorakṣa, Kāki, Kāpālika and other [perfected yogis].\endnote{It is not clear who Kāki and Kāpālika might be, or whether kākikāpālika was the intended name. In fact, it seems somewhat likely that the compound \textit{°kākikāpālika°} is a corruption of \textit{°khaṇḍakāpālika°}, the name of a \textit{siddha} mentioned in some manuscripts of the \textit{Haṭhapradīpikā} (1.8).} \textbf{Celibacy} is when not [even] a drop of semen falls through the opening of the penis, and female generative fluid and [other substances] are drawn up [through the urethra]. \textbf{Silence} is avoiding talk about worldly matters. \textbf{Steadiness} is the stillness of mind [brought about] by not worrying about caring for a wife, son and [family]. \textbf{Patience} is that forbearance [required] when physical and mental suffering is caused by thieves, kings and bad people and when one has the power to punish them. \textbf{Fearlessness} is not being afraid as a result of knowing that the body as the locus of ego and one's sons and houses as the seat for the sense of ownership, etc., are impermanent, and that the self, which is truth, consciousness and bliss, is permanent.

\textbf{Purity} is purification of the body by [cleansing techniques], such as enemas, swallowing a strip of cloth, and emesis; purification of speech by singing the names of god; and purification of mind by meditating on god. 

\textbf{Mantra recitation} is the regular repetition of the mantra of the deity to be worshipped. \textbf{Asceticism} is carrying out one's religious duty (\textit{svadharma}). For the householder, celibate student and forest hermit, a \textbf{fire sacrifice} is the daily fire sacrifice;  for ascetics, Brahmins only in name, women and the like, it is the fire sacrifice that is an ancillary within the worship of god in [scriptures, such as] Vaiṣṇava Tantras. \textbf{Confidence} is eagerness for following the rules taught in scripture. \textbf{Hospitality} is kindness towards those who have practised the methods of one’s own path, such as Haṭhayoga. [\textbf{Worship of me}] is worshipping god, [that is,] compassion for all beings, complete contentment with anything, and seeing the self in all things. \textbf{Travelling to sacred sites} is going to various places for the sake of seeing those who have perfected the practice [of yoga].  \textbf{Striving to help others} is effort directed at [helping] students succeed at the practice.  \textbf{Satisfaction} is contentment with what has been obtained by chance. \textbf{Serving the teacher} is procuring the favour of a teacher who is deeply versed in scripture and the experience [of yoga].
\bigskip

\Section[Postures (\textit{āsana})]{[Postures (\textit{āsana})]}

Next, the postures for attaining the capacity for [practising] the six cleansing techniques are described.
\end{eng}
\end{multicols}


\newpage
~\vskip 1cm

\begin{center}
\begin{kan}
\booktitle{ಕಪಾಲ ಕುರಂಟಕನ ಹಠಯೋಗ ಪದ್ಧತಿಯ ಕೈಪಿಡಿ} 
%\addcontentsline{tocka}{chapter}{\kantoc{ಕಪಾಲ ಕುರಂಟಕನ ಹಠಯೋಗ ಪದ್ಧತಿಯ ಕೈಪಿಡಿ}}
\medskip

{\Large\bfseries ಮಹಾಗಣಪತಯೇ ನಮಃ  ಶ್ರೀ ಗುರುಭ್ಯೋ ನಮಃ}

{\Large\bfseries ಮಹಾಗಣಪತಿ ಹಾಗು ಪೂಜ್ಯ ಗುರುಗಳಿಗೆ ನಮನಗಳು}
\end{kan}
\end{center}
\bigskip


\begin{multicols}{2}
\begin{kan}
ಈ ಹಠಯೋಗ ಅಭ್ಯಾಸ ಕೈಪಿಡಿಯನ್ನು ಕಪಾಲಕುರಂಟಕರು ರಚಿಸಿದ್ದಾರೆ. ಇದನ್ನು ವಿಶೇಷವಾಗಿ ಜನನ-ಮರಣಗಳ ಚಕ್ರದಿಂದ ಪೀಡಿತ\-ರಾದವರು, ಪಂಚೇಂದ್ರಿಯಗಳ ಸುಖಕ್ಕೆ ಅತಿಯಾಗಿ ಅಂಟಿಕೊಂಡ\-ವರು, ಸ್ತ್ರೀ ವ್ಯಾಮೋಹಕ್ಕೊಳಗಾದವರು, ಸಮಾಜದಿಂದ ಬಹಿಷ್ಕೃತರಾದವರು ಹಾಗು ಪಾಪ ಕಾರ್ಯಗಳನ್ನು ಮಾಡುವವರಿಗಾಗಿ ರಚಿಸಲಾಗಿದೆ. ಈ ಗ್ರಂಥದಲ್ಲಿ ಹಠಯೋಗದ ವಿವಿಧ ವಿಷಯಗಳು ಮತ್ತು ಅವುಗಳನ್ನು ಅಭ್ಯಾಸ ಮಾಡುವ ವಿಧಾನಗಳ ಕುರಿತು ವಿವರವಾಗಿ ಬರೆಯಲಾಗಿದೆ.
\bigskip

\Section{(ಯೋಗಿಯ) ಕುಟೀರದ ಲಕ್ಷಣಗಳು}

ಕುಟೀರವು ನಾಲ್ಕು ಮುಂಗೈಯಷ್ಟು ಎತ್ತರ ಮತ್ತು ಅಷ್ಟೇ ಅಗಲವಿರಬೇಕು. ಯೋಗಮುದ್ರಾಭ್ಯಾಸಕ್ಕಾಗಿ ಕುಟೀರವನ್ನು ಬೂದಿಯಿಂದ ಲೇಪಿಸಿರಬೇಕು. ಆಸನಗಳ ಅಭ್ಯಾಸಕ್ಕಾಗಿ ಕುಟೀರವನ್ನು ಕೆಮ್ಮಣಿನಿಂದ ಲೇಪಿಸಿರಬೇಕು. ಬಸ್ತಿಯಂತಹ ಆರು ಶೋಧನಾ ಕ್ರಿಯೆಗಳ ಅಭ್ಯಾಸಕ್ಕಾಗಿ ಕುಟೀರವನ್ನು ಸುಣ್ಣದಿಂದ ಸಾರಿಸಿರಬೇಕು. ವಿಶ್ರಾಂತಿಗಾಗಿ ಕುಟೀರದಲ್ಲಿ ಹುಲಿ ಅಥವಾ ಇತರ ಪ್ರಾಣಿಗಳ ಚರ್ಮವನ್ನು ಹಾಸಿರಬೇಕು. ವಜ್ರೊಲಿ ಕ್ರಿಯೆ ಅಭ್ಯಾಸಕ್ಕಾಗಿ ಹತ್ತಿಯ ಬಟ್ಟೆಯನ್ನು ಬಳಸಬೇಕು. ಹಗ್ಗದ ಸಹಾಯದಿಂದ ಮಾಡುವ ಅಥವಾ  ಎತ್ತರಕ್ಕೆ ಏರುವ ಆಸನಗಳ (ಸ್ವರ್ಗಾರೋಹಣ) ಅಭ್ಯಾಸಕ್ಕಾಗಿ ಕುಟೀರವು ಮೂರು ಬಿಲ್ಲುಗಳಷ್ಟು ಎತ್ತರ ಮತ್ತು ಒಂದು ಬಿಲ್ಲಿನಷ್ಟು ಅಗಲವಿರಬೇಕು. ಮಲ ಮೂತ್ರ. ವಿಸರ್ಜನೆಗಾಗಿ ಹತ್ತಿರದಲ್ಲಿ ಗುಪ್ತ ಸ್ಥಳವಿರಬೇಕು. ಸ್ನಾನಕ್ಕಾಗಿ ನಿರ್ಜನ ಪ್ರದೇಶದಲ್ಲಿ ನಿರ್ಮಲವಾದ ಹರಿಯುವ ನೀರಿರಬೇಕು.

\vphantom{111}\\

\Section{ಯಮ ಮತ್ತು ನಿಯಮಗಳ ಲಕ್ಷಣಗಳು}

ಅಹಿಂಸೆ, ಸತ್ಯಪಾಲನೆ, ಕಳ್ಳತನ ಮಾಡದಿರುವುದು, ವಾತ್ಸಲ್ಯವಿಲ್ಲದಿರುವುದು, ಲಜ್ಜೆ , ಅಸಂಚಯ, ಆಸ್ತಿಕತೆ,  ಬ್ರಹ್ಮಚರ್ಯ, ಮೌನ, ಸ್ಥಿರಮನಸ್ಥಿತಿ, ತಾಳ್ಮೆ,ಮತ್ತು ಧೈರ್ಯ ಇವುಗಳನ್ನು ‘ಯಮ’ ಎಂದು ಕರೆಯಲಾಗುತ್ತದೆ.  ಶುಚಿತ್ವ, ಮಂತ್ರ ಜಪ, ತಪಸ್ಸು, ಹೋಮ, ಶ್ರದ್ಧೆ, ಆತಿಥ್ಯ, ದೈವಾರಾಧನೆ, ತೀರ್ಥಯಾತ್ರೆ, ಪರೋಪಕಾರ, ತೃಪ್ತಿ ಮತ್ತು ಗುರುಸೇವೆ ಇವುಗಳನ್ನು ‘ನಿಯಮ’ ಎಂದು ಹೇಳಲಾಗುತ್ತದೆ.

ಕಾಯಾ ವಾಚಾ ಮನಸಾ ಎಂದರೇ, ಇತರರ ದೇಹಕ್ಕೆ, ಮಾತಿನಿಂದ ಅಥವಾ ಮನಸ್ಸಿಗೆ ನೋವನ್ನುಂಟು ಮಾಡದಿರುವುದೇ \textbf{ಅಹಿಂಸೆ}. ಯಾವುದು ವಾಸ್ತವವೋ, ಯಾವುದು ಇತರರಿಗೆ ಹಾನಿ ಮಾಡುವುದಿಲ್ಲವೋ, ಮನಸ್ಸನ್ನು ಅಚಲವಾಗಿಸುವುದೋ, ಯಾವ ಮಾತು ಹಿತವನ್ನು ನೀಡುತ್ತದೆಯೋ ಅಂತಹ ಮಾತೇ \textbf{ಸತ್ಯ}. ಪರರ ವಸ್ತುಗಳ ಮೇಲೆ ಮನಸ್ಸು, ಮಾತು ಅಥವಾ ಕ್ರಿಯೆಯ ಮೂಲಕ ಆಸೆ ಪಡದಿರುವುದೇ \textbf{ಕಳ್ಳತನ ಮಾಡದಿರುವುದು}. ಯೋಗಾಭ್ಯಾಸವನ್ನು ವಿರೋಧಿಸುವವರ ಮೇಲೆ ಸ್ನೇಹಭಾವವಿಲ್ಲದಿರುವುದೇ \textbf{ವಾತ್ಸಲ್ಯವಿಲ್ಲದಿರುವುದು}. ಧರ್ಮ ಅಥವಾ ಅಧರ್ಮ ಮಾರ್ಗದಿಂದ ಒಬ್ಬರ ಯೋಗಾಭ್ಯಾಸಕ್ಕೆ ಧಕ್ಕೆಯುಂಟು ಮಾಡಿ ಅದಕ್ಕೆ ಪಶ್ಚಾತಾಪ ಪಡುವುದೇ \textbf{ಲಜ್ಜೆ}. ಯೋಗಾಭ್ಯಾಸಕ್ಕೆ ಬೇಕಾಗುವ ವಸ್ತುಗಳ ಹೊರತಾಗಿ ಇತರ ವಸ್ತುಗಳನ್ನು ಅತಿಯಾಗಿ ಸಂಗ್ರಹಮಾಡದಿರುವುದು \textbf{ಅಸಂಚಯ}. ಮತ್ಸ್ಯೇಂದ್ರನಾಥರು, ಗೋರಕ್ಷನಾಥರು, ಕಾಕಿ, ಕಾಪಾಲಿಕರು ಮುಂತಾದ ಸಿದ್ಧಯೋಗಿಗಳ ಶಾಸ್ತ್ರಗಳಲ್ಲಿ ಅಚಲ ನಂಬಿಕೆಯಿರುವುದೇ \textbf{ಆಸ್ತಿಕತೆ}. ಶಿಶ್ನ ರಂದ್ರದಿಂದ ಒಂದು ತೊಟ್ಟೂ ವೀರ್ಯ ನಷ್ಟವಾಗದಂತೆ ಹಾಗು ಸ್ತ್ರೀಯು ತನ್ನ ಉತ್ಪಾದಕ ದ್ರವವನ್ನು ಯೋನಿಯಲ್ಲೇ  ಹಿಡಿದಿಟ್ಟುಕೊಳ್ಳುವ ಶಕ್ತಿಯೇ \textbf{ಬ್ರಹ್ಮಚರ್ಯ}. ಲೌಕಿಕ ವಾದವಿವಾದಗಳನ್ನು ಬಿಡುವುದೇ \textbf{ಮೌನ}. ಹೆಂಡತಿ, ಮಕ್ಕಳು ಮತ್ತು ಸಂಸಾರದ ಚಿಂತೆಯಿಲ್ಲದೇ ಮನಸ್ಸನ್ನು ಅಚಲವಾಗಿಟ್ಟುಕೂಳ್ಳುವುದೇ \textbf{ಸ್ಥಿರಮನಸ್ಥಿತಿ}. ಕಳ್ಳರು, ರಾಜರು, ದುಷ್ಟರಿಂದ ದೈಹಿಕ ಅಥವಾ ಮಾನಸಿಕವಾಗಿ ತೊಂದರೆಯಾದರೂ, ಅವರನ್ನು ಶಿಕ್ಷಿಸುವ ಶಕ್ತಿಯಿದ್ದರೂ ಸಹಿಸಿಕೊಳ್ಳುವುದೇ \textbf{ತಾಳ್ಮೆ}. ದೇಹ, ಅಹಂಕಾರ, ಸಂಸಾರ, ಆಸ್ತಿಗಳೆಲ್ಲವೂ ತಾತ್ಕಾಲಿಕ. ಆದರೆ, ಸಚ್ಚಿದಾನಂದ ರೂಪವಾದ ಆತ್ಮವು ಮಾತ್ರ ಶಾಶ್ವತ ಎಂಬುದನ್ನರಿತು ಭಯವಿಲ್ಲದೇ ಇರುವುದೇ  \textbf{ಧೈರ್ಯ}.

ಬಸ್ತಿ ಮತ್ತು ಧೌತಿಗಳಂತಹ ಕ್ರಿಯೆಗಳಿಂದ ದೇಹವನ್ನು, ದೇವರ ನಾಮಸ್ಮರಣೆಯಿಂದ ಮಾತನ್ನು ಮತ್ತು ಧ್ಯಾನದಿಂದ ಮನಸ್ಸನ್ನು ಪವಿತ್ರಗೊಳಿಸುವುದೇ \textbf{ಶುಚಿತ್ವ}. ದೇವರ ಮಂತ್ರವನ್ನು ನಿರಂತರವಾಗಿ ಜಪಿಸುವುದೇ \textbf{ಮಂತ್ರ ಜಪ}. ತನ್ನ ಸ್ವಧರ್ಮವನ್ನು ಪಾಲಿಸುವುದೇ \textbf{ತಪಸ್ಸು}. ಗೃಹಸ್ಥರು, ಬ್ರಹ್ಮಚಾರಿಗಳು, ವಿದ್ಯಾರ್ಥಿಗಳು ಮತ್ತು ವನವಾಸಿಗಳಿಗೆ ನಿತ್ಯ ಅಗ್ನಿಹೋತ್ರವು; ಬ್ರಾಹ್ಮಣರು, ಮಹಿಳೆಯರು ಮತ್ತು ಇತರರಿಗೆ ವೈಷ್ಣವ ತಂತ್ರದಂತಹ ಶಾಸ್ತ್ರಗಳ ಪ್ರಕಾರ ದೈವಾರಾಧನೆಯ ಅಂಗವಾಗಿ ಮಾಡುವ ಯಜ್ಞವೇ \textbf{ಹೋಮ}. ಶಾಸ್ತ್ರದ ನಿಯಮಗಳನ್ನು ನಿಷ್ಠೆಯಿಂದ ಪಾಲಿಸುವ ಉತ್ಸಾಹವೇ \textbf{ಶ್ರದ್ಧೆ}. ತನ್ನ ಹಾದಿಯಲ್ಲೇ ಹಠಯೋಗದಂತಹ ಸಾಧನೆ ಮಾಡುವ ಯೋಗಿಗಳಿಗೆ ತೋರುವ ಪ್ರೀತಿಯೇ \textbf{ಆತಿಥ್ಯ}. ಎಲ್ಲಾ ಜೀವಿಗಳ ಮೇಲೆ ಕರುಣೆ, ಎಲ್ಲದರಲ್ಲೂ ತೃಪ್ತಿ ಮತ್ತು ಎಲ್ಲದರಲ್ಲೂ ತನ್ನನ್ನು ಕಾಣುವುದೇ \textbf{ದೈವಾರಾಧನೆ}. ಸಿದ್ಧ ಪುರುಷರನ್ನು ನೋಡಲು ವಿವಿಧ ಪುಣ್ಯಕ್ಷೇತ್ರಗಳಿಗೆ ಪ್ರಯಾಣಿಸುವುದೇ \textbf{ತೀರ್ಥಯಾತ್ರೆ}.  ಶಿಷ್ಯರಿಗೆ ಯೋಗ ಸಾಧನೆಯಲ್ಲಿ ಸಹಕರಿಸುವುದು ಹಾಗು ಸಹಾಯ ಮಾಡುವುದೇ \textbf{ಪರೋಪಕಾರ}. ಅನಿರೀಕ್ಷಿತವಾಗಿ ಮತ್ತು ಆಕಸ್ಮಿಕವಾಗಿ ಸಿಕ್ಕಿದುದರಲ್ಲೇ ಸಂತೋಷಪಡುವುದು \textbf{ತೃಪ್ತಿ}. ಯೋಗದ ಅನುಭವ ಮತ್ತು ಶಾಸ್ತ್ರಜ್ಞಾನವಿರುವ ಗುರುವಿನ ಪ್ರೀತಿ ಮತ್ತು ಅನುಗ್ರಹಕ್ಕೆ ಪಾತ್ರರಾಗುವುದೇ \textbf{ಗುರುಸೇವೆ}.
\bigskip

\Section{ಆಸನಗಳು}

ಮುಂದೆ, ಆರು ಶುದ್ಧಿಕ್ರಿಯೆಗಳನ್ನು ಅಭ್ಯಾಸ ಮಾಡುವ ಸಾಮಥ್ರ್ಯವನ್ನು ಪಡೆಯಲು ಸಹಕಾರಿಯಾದ ಆಸನಗಳನ್ನು ವಿವರಿಸಲಾಗಿದೆ. 
\end{kan}
\end{multicols}
